\chapter*{General Topological Spaces}
\phantomsection
\addcontentsline{toc}{chapter}{General Topological Spaces}
\markboth{General Topological Spaces}{General Topological Spaces}

The metric chapters developed open sets, continuity, compactness, and connectedness by starting from a distance function. The same ideas make sense more generally when the family of open sets is taken as the primitive structure. This interlude records that abstraction explicitly, so that the phrase \emph{topological foundations} refers to more than metric topology while keeping the numbered II/01--II/25 architecture unchanged.

\section{Topologies: open sets as primitive data}
Let \(X\) be a set. A \emph{topology} \(\mathcal T\) on \(X\) is a collection of subsets of \(X\) satisfying
\[
\varnothing,X\in\mathcal T,\qquad
\bigcup_{\alpha\in A}U_\alpha\in\mathcal T,
\qquad
U_1\cap\cdots\cap U_m\in\mathcal T
\]
for every family \(\{U_\alpha\}_{\alpha\in A}\subseteq\mathcal T\) and every finite family \(U_1,\dots,U_m\in\mathcal T\). The members of \(\mathcal T\) are the \emph{open sets}; complements of open sets are \emph{closed}.

Every metric \(d\) produces a topology: a set is open when each of its points contains a metric ball lying inside the set. Once that topology has been extracted, many statements no longer need the numerical values of distances.

\begin{example}[Three topologies on one set]
On a set \(X\), the discrete topology \(\mathcal P(X)\) declares every subset open, while the indiscrete topology \(\{\varnothing,X\}\) declares only the two forced sets open. On \(\mathbb R\), the usual topology lies between these extremes and is generated by open intervals.
\end{example}

A property that depends only on the topology, and not on a particular compatible metric, is called \emph{topological}. Continuity, compactness, connectedness, and homeomorphism type are topological; a numerical distance, angle, or Lipschitz constant is not.

\section{Bases and subbases}
A family \(\mathcal B\subseteq\mathcal T\) is a \emph{basis} if every open set is a union of members of \(\mathcal B\). Equivalently, for every open \(U\) and every \(x\in U\), there is \(B\in\mathcal B\) with
\[
x\in B\subseteq U.
\]
For the usual topology of \(\mathbb R\), open intervals form a basis. Intervals with rational endpoints already form a countable basis.

A family \(\mathcal S\) is a \emph{subbasis} when finite intersections of members of \(\mathcal S\) form a basis. Subbases are useful because they let a topology be specified from a small collection of elementary conditions.

\paragraph{Basis criterion.}
Suppose \(\mathcal B\) covers \(X\), and whenever \(x\in B_1\cap B_2\) with \(B_1,B_2\in\mathcal B\), there is \(B_3\in\mathcal B\) such that
\[
x\in B_3\subseteq B_1\cap B_2.
\]
Then arbitrary unions of elements of \(\mathcal B\) form a topology. The local refinement condition is exactly what is needed to make finite intersections open.

\section{Subspace topology}
If \((X,\mathcal T)\) is a topological space and \(Y\subseteq X\), the \emph{subspace topology} on \(Y\) is
\[
\mathcal T_Y=\{Y\cap U:U\in\mathcal T\}.
\]
Thus openness is always relative to the space under discussion. For example, \([0,1)\) is not open in \(\mathbb R\), but it is open in the subspace \([0,2]\) because
\[
[0,1)=(-1,1)\cap[0,2].
\]

If \(F\subseteq Y\), then \(F\) is closed in \(Y\) exactly when \(F=Y\cap C\) for some closed \(C\subseteq X\). This is the closed-set counterpart of the same construction.

\section{Finite products}
For topological spaces \(X\) and \(Y\), the product topology on \(X\times Y\) has basis
\[
\{U\times V:U\subseteq X\text{ open},\;V\subseteq Y\text{ open}\}.
\]
For a finite product \(X_1\times\cdots\times X_n\), basic open sets are products \(U_1\times\cdots\times U_n\) of open coordinate sets.

On \(\mathbb R^n\), this product topology agrees with the usual Euclidean topology. One way to see the equivalence is that every Euclidean ball contains a sufficiently small open box about its center, while every open box contains a sufficiently small Euclidean ball.

\begin{example}[Coordinate projections]
The projection \(\pi_X:X\times Y\to X\), \(\pi_X(x,y)=x\), is continuous because
\[
\pi_X^{-1}(U)=U\times Y
\]
is open whenever \(U\) is open. Product topology is the natural topology making the coordinate projections continuous.
\end{example}

\section{Continuity by inverse images}
For topological spaces \(X,Y\), a map \(f:X\to Y\) is \emph{continuous} if
\[
V\subseteq Y\text{ open}\quad\Longrightarrow\quad f^{-1}(V)\subseteq X\text{ open}.
\]
This definition contains no metric and is equivalent to the usual epsilon--delta definition when both spaces are metric.

Inverse images are the right operation because they preserve unions, finite intersections, and complements:
\[
f^{-1}\!\left(\bigcup_\alpha V_\alpha\right)=\bigcup_\alpha f^{-1}(V_\alpha),
\qquad
f^{-1}(V\cap W)=f^{-1}(V)\cap f^{-1}(W).
\]
Consequently it is enough to test continuity on a basis, or even on a subbasis, for the topology of the target.

\paragraph{Composition.}
If \(f:X\to Y\) and \(g:Y\to Z\) are continuous, then \(g\circ f\) is continuous, because
\[
(g\circ f)^{-1}(U)=f^{-1}(g^{-1}(U)).
\]
The proof is purely topological.

\section{Homeomorphisms and topological invariants}
A bijection \(f:X\to Y\) is a \emph{homeomorphism} if both \(f\) and \(f^{-1}\) are continuous. Homeomorphic spaces are regarded as the same from the viewpoint of topology.

The distinction between a continuous bijection and a homeomorphism matters: continuity controls inverse images in one direction, while a homeomorphism requires the topology to be transported exactly in both directions. A standard useful criterion is that a continuous bijection from a compact space to a Hausdorff space is automatically a homeomorphism.

\paragraph{Metric versus topological information.}
A homeomorphism need not preserve distances. The map
\[
f:\mathbb R\to(-1,1),\qquad f(x)=\frac{x}{1+|x|},
\]
is a homeomorphism, although the two spaces have very different metric sizes. Boundedness with respect to a chosen metric is therefore not, by itself, a topological invariant.

\section{Compactness beyond metrics}
A topological space \(X\) is \emph{compact} if every open cover of \(X\) has a finite subcover. This is the general definition; no sequence or distance is required.

\paragraph{Continuous images preserve compactness.}
If \(X\) is compact and \(f:X\to Y\) is continuous, then \(f(X)\) is compact. Indeed, an open cover of \(f(X)\) pulls back to an open cover of \(X\); choose finitely many pulled-back sets and then return to the corresponding original sets.

In metric spaces, compactness is equivalent to sequential compactness. In arbitrary topological spaces that equivalence can fail. The open-cover definition is therefore the correct general starting point.

\section{Connectedness beyond metrics}
A topological space \(X\) is \emph{connected} if it cannot be written as a union
\[
X=U\cup V
\]
of two disjoint nonempty open sets. Equivalently, the only subsets that are both open and closed are \(\varnothing\) and \(X\).

\paragraph{Continuous images preserve connectedness.}
If \(X\) is connected and \(f:X\to Y\) is continuous, then \(f(X)\) is connected. A separation of \(f(X)\) would pull back to a separation of \(X\).

Path connectedness still implies connectedness in every topological space, but the converse need not hold. Thus the distinction made in Chapter~7 is genuinely topological rather than specifically metric.

\section{Hausdorff separation and uniqueness of limits}
A topological space is \emph{Hausdorff} if any two distinct points have disjoint neighborhoods. Metric spaces are Hausdorff: if \(x\ne y\), balls of radius smaller than \(d(x,y)/2\) separate them.

Hausdorff separation explains why limits, when sequences are an adequate tool, are unique. Without a separation axiom, a sequence may converge to more than one point. This shows that some familiar facts from metric analysis are consequences of additional topological structure, not consequences of the topology axioms alone.

\section{Why sequences are not always enough}
Metric spaces are \emph{first countable}: at each point \(x\), the balls
\[
B(x,1),\ B(x,1/2),\ B(x,1/3),\dots
\]
form a countable local basis. This is why sequences characterize closure and continuity so effectively in Chapters~3--7.

In an arbitrary topological space, a point may have no countable local basis. Then the statement
\[
x\in\overline A
\]
need not imply that some sequence of points of \(A\) converges to \(x\). Likewise, sequential tests for continuity or compactness may cease to be equivalent to their open-set formulations.

The practical rule for this volume is therefore:
\begin{quote}
Sequence arguments are fully justified in the metric spaces developed in Chapters~2--7. When a later argument is phrased for an arbitrary topological space, use open sets, neighborhoods, covers, and inverse images unless an additional countability hypothesis has been stated.
\end{quote}

\section{A compact dictionary}
The following table separates the metric realization from the underlying topological idea.

\begin{center}
\begin{tabular}{p{0.27\textwidth}p{0.31\textwidth}p{0.31\textwidth}}
\textbf{Idea} & \textbf{Metric formulation} & \textbf{Topological formulation}\\\hline
Neighborhood & contains a ball about the point & contains an open set about the point\\
Continuity & epsilon--delta control & inverse images of open sets are open\\
Compactness & sequential/open-cover equivalents & finite subcover property\\
Connectedness & no metric separation needed & no separation into two nonempty open sets\\
Equivalence of spaces & isometry is strongest & homeomorphism preserves topology\\
Closure test & sequences suffice & neighborhoods always suffice\\
\end{tabular}
\end{center}

\section{Scope boundary}
This interlude supplies the general language promised by the volume title, but it is deliberately not a course in algebraic topology. In particular, the higher-dimensional Brouwer fixed-point theorem in Chapter~17 still requires additional machinery---for example Sperner's lemma, no-retraction methods, degree theory, or homology---that is not developed here.

The distinction is important. General topology provides the vocabulary of continuity, compactness, connectedness, products, and homeomorphisms. Brouwer's theorem is a deeper global result about Euclidean balls and compact convex sets. Stating the general definitions does not by itself constitute a proof of Brouwer.

\section*{Worked checks}
\begin{enumerate}
\item On a discrete space, every map to any topological space is continuous because every inverse image is open.
\item A map into an indiscrete space is always continuous because the only open target sets are \(\varnothing\) and the whole target.
\item If \(Y\subseteq X\) has the subspace topology, the inclusion \(Y\hookrightarrow X\) is continuous.
\item In a product \(X\times Y\), both coordinate projections are continuous.
\item A continuous image of a compact connected space is again compact and connected.
\item In a Hausdorff space, a compact subset is closed; this is one of the main reasons Hausdorff spaces behave like familiar metric spaces.
\end{enumerate}

\section*{Exercises}
\begin{enumerate}
\item Verify directly that the discrete and indiscrete collections are topologies on every set.
\item Prove that arbitrary intersections of closed sets and finite unions of closed sets are closed.
\item Show that \([0,1)\) is open in \([0,2]\) with the subspace topology and determine whether it is closed there.
\item Prove that a map \(f:X\to Y\) is continuous if the inverse image of every member of a basis for \(Y\) is open in \(X\).
\item Show that the product of two connected spaces is connected.
\item Prove that a continuous bijection from a compact space to a Hausdorff space is a homeomorphism.
\item Give two different metrics on \(\mathbb R\) that induce the same topology, and identify a metric quantity they do not preserve.
\item Explain precisely where first countability enters the metric proof that \(x\in\overline A\) implies the existence of a sequence \((a_n)\subseteq A\) converging to \(x\).
\item Decide which of the following are topological properties: boundedness, compactness, connectedness, completeness of a chosen metric, Hausdorffness, and path connectedness.
\item Explain why the material in this interlude clarifies the statement of Brouwer's theorem but does not replace a proof of it.
\end{enumerate}

\section*{Interlude summary}
A topology abstracts the open-set structure produced by a metric. Bases and subbases generate that structure efficiently; subspaces and finite products construct new spaces; continuity is expressed by inverse images; homeomorphisms identify spaces with the same topology; compactness and connectedness extend without distances; and first countability explains why sequence-based reasoning is especially effective in metric analysis. These ideas provide the general topological foundation needed for later volumes while keeping deeper fixed-point and algebraic-topological machinery in its proper place.
